%% ch06_microgrid_protection.tex
%% Chapter 6: Protection of IBR Microgrids
%% ============================================================
\chapter{Protection of Inverter-Based Microgrids}
\label{ch:microgrid}

\section{Introduction}
\label{sec:ch5mg_intro}

A microgrid is a controllable cluster of generation (predominantly IBR),
storage, and load that can operate either connected to the main grid or
in an islanded autonomous mode. The transition between these operating modes
is the most safety-critical event in microgrid operation: the protection
must reconfigure within milliseconds to handle the radical change in fault
current magnitude and direction.

This chapter delivers three contributions:
\begin{itemize}
\item GFL/GFM dual-mode characterisation for protection
      coordination (Section~\ref{sec:ch5_characterisation}).
\item ROCOF-supervised islanding detection with IEC~61850
      mode-switch achieving isolation within 120~ms
      (Section~\ref{sec:ch5_islanding}).
\item Adaptive OC setting-group architecture triggered by
      GOOSE mode flags (Section~\ref{sec:ch5_adaptiveoc}).
\end{itemize}

\section{GFL vs.\ GFM: Protection-Relevant Characterisation}
\label{sec:ch5_characterisation}

\subsection{Operating Mode Definition}

\begin{thesisdef}{Grid-Connected Mode}
The microgrid is synchronised to the main grid. The grid provides the
voltage reference and limits frequency excursions. IBRs operate in GFL
mode (PLL-locked current source) unless VSM/GFM firmware is engaged.
Fault current levels are high because the grid infeed dominates.
\end{thesisdef}

\begin{thesisdef}{Islanded Mode}
The microgrid operates autonomously without grid connection. All voltage
and frequency regulation must be provided by internal IBRs. Only GFM
inverters (or droop-controlled GFL with synthetic inertia) can maintain
stable islanded operation. Fault current is limited to $I_\mathrm{lim,GFM}
\approx 1.5$--$2.5$~pu from GFM units.
\end{thesisdef}

\subsection{Fault Current Comparison}

Table~\ref{tab:ch6_fault_comparison} compares the fault current levels
in grid-connected versus islanded mode for the SAMBP microgrid test system.

\begin{table}[htbp]
\caption{Fault current levels in microgrid operating modes (33~kV, 10~MVA base)}
\label{tab:ch6_fault_comparison}
\begin{tabularx}{\textwidth}{@{}lXXX@{}}
\toprule
Fault type & Grid-connected (pu) & Islanded GFL (pu) & Islanded GFM (pu) \\
\midrule
3-phase (3P)       & 6.0--10.0 & 1.1--1.5 & 1.5--2.5 \\
Single LG (SLG)    & 4.0--8.0  & 0.7--1.0 & 1.0--2.0 \\
Line-to-line (LL)  & 5.0--9.0  & 1.0--1.3 & 1.3--2.2 \\
High-impedance (HIF) & 0.1--0.5 & 0.05--0.2 & 0.1--0.3 \\
\bottomrule
\end{tabularx}
\end{table}

The fault current reduction from grid-connected to islanded is 4--8$\times$
for GFL-dominated islands and 3--5$\times$ for GFM-dominated islands.
This reduction is too large for the protection to accommodate with a
single setting group.

\subsection{GFL/GFM Transition Dynamics}

When the main breaker opens (island detection), GFL inverters must
transition to GFM mode within $T_\mathrm{switch} \leq 100$~ms to
maintain island voltage stability. The transition involves:
\begin{enumerate}
\item Disabling the PLL (GFL synchronisation reference lost).
\item Activating the VSM outer loop with virtual inertia $M$ and damping $D$.
\item Enabling reactive power droop: $Q = Q_0 - k_V(|V| - |V_0|)$.
\end{enumerate}
During the transition window $[0, T_\mathrm{switch}]$, the IBR is in
an undefined state—neither GFL nor GFM. The protection system must
withstand this transition without false tripping.

The protection challenge is that during this window:
\begin{itemize}
\item $\kibr$ changes rapidly from $\kibr^\mathrm{GFL} \leq 1.5$~pu
      to $\kibr^\mathrm{GFM} \approx 2.0$~pu.
\item The current angle $\phi$ jumps by up to $90^\circ$ as the reactive
      injection priority changes.
\item The frequency deviates by up to $\pm 2$~Hz before the VSM
      governor action stabilises it.
\end{itemize}

\section{ROCOF-Supervised Islanding Detection}
\label{sec:ch5_islanding}

\subsection{Passive Islanding Detection Methods}

IEEE~1547-2018 mandates that all IBRs cease energizing the grid within
2~s of islanding. Passive islanding detection methods monitor local
electrical quantities:
\begin{itemize}
\item \textbf{Over/Under Voltage (OV/UV):} Detects islanding when
      $V < 0.88$~pu or $V > 1.10$~pu. Fails when load exactly equals
      generation (near-unity load balance).
\item \textbf{Over/Under Frequency (OF/UF):} Detects islanding when
      $f < 59.3$~Hz or $f > 60.5$~Hz (IEEE~1547). Fails at load balance.
\item \textbf{ROCOF:} Rate-of-Change-of-Frequency. Detects the
      characteristic frequency ramp caused by power imbalance at the
      moment of islanding.
\end{itemize}

\subsection{ROCOF Theory}

The frequency dynamics after islanding are governed by the combined
swing equation:
\begin{equation}
  M_\mathrm{total} \frac{d^2\theta}{dt^2} = P_g - P_l
  \label{eq:ch6_swing}
\end{equation}
where $M_\mathrm{total} = \sum_i M_i$ is the total virtual inertia of
all GFM units, $P_g$ is total generation, and $P_l$ is total load.
The ROCOF immediately after islanding is:
\begin{equation}
  \left.\frac{df}{dt}\right|_{t=0^+} = \frac{f_0 (P_g - P_l)}{2H_\mathrm{total} S_\mathrm{rated}}
  \label{eq:ch6_rocof}
\end{equation}
where $H_\mathrm{total}$ is the total virtual inertia constant (seconds)
and $S_\mathrm{rated}$ is the aggregate rated power. For a 10\% load-generation
imbalance with $H = 2$~s, $S = 10$~MVA:
\begin{equation}
  \left.\frac{df}{dt}\right|_{t=0^+}
  = \frac{50 \cdot 0.10 \cdot 10}{2 \cdot 2 \cdot 10} = 1.25\;\mathrm{Hz/s}
  \label{eq:ch6_rocof_example}
\end{equation}

\subsection{ROCOF Calculation}

The ROCOF is estimated by fitting a linear frequency trend over a
sliding window of $N_w$ cycles:
\begin{equation}
  \widehat{\mathrm{ROCOF}}(t) = \frac{f(t) - f(t - N_w/f_0)}{N_w/f_0}
  \label{eq:ch6_rocof_est}
\end{equation}
A window $N_w = 3$ cycles (60~ms at 50~Hz) provides a good trade-off
between speed and noise rejection.

\subsection{Vector Surge Supplement}

The vector surge element detects the instantaneous phase angle jump
$\Delta\theta$ at the moment of islanding:
\begin{equation}
  \Delta\theta = \arctan\!\left(\frac{\omega_0 P_\mathrm{imbalance}}
    {2 H_\mathrm{total} S_\mathrm{rated} (df/dt)}\right)
  \label{eq:ch6_vector_surge}
\end{equation}
For $P_\mathrm{imbalance} = 10\%$: $\Delta\theta \approx 3$--$8^\circ{}$ (too
small to detect reliably with classical $2^\circ$ threshold). The vector surge
element is used only as confirmation, not primary detection.

\subsection{Combined ROCOF + Impedance Islanding Criterion}

The proposed islanding detection criterion combines ROCOF, vector surge,
and negative-sequence voltage ratio:
\begin{equation}
  T_\mathrm{island} = \mathbf{1}\!\left[|df/dt| > \gamma_1\right]
  \;\cap\; \mathbf{1}\!\left[|\Delta\theta| > \gamma_2\right]
  \;\cup\; \mathbf{1}\!\left[V_2/V_1 > \gamma_3\right]
  \label{eq:ch6_island_logic}
\end{equation}
where:
\begin{itemize}
\item $\gamma_1 = 0.5$~Hz/s (ROCOF threshold; below inadvertent ROCOF
      for grid disturbances).
\item $\gamma_2 = 2^\circ{}$ (vector surge threshold).
\item $\gamma_3 = 0.02$ (negative-sequence voltage ratio threshold;
      detects unbalanced load after islanding).
\end{itemize}

\textbf{Non-Detection Zone (NDZ) Analysis:}
The NDZ is the set of load-generation balance conditions where the
islanding detector fails. For the ROCOF element:
\begin{equation}
  \mathrm{NDZ}_\mathrm{ROCOF} = \left\{(P_g, P_l):\;
    \left|\frac{f_0(P_g - P_l)}{2H S_r}\right| < \gamma_1 \right\}
  = \left\{P_\mathrm{imb} < \frac{2H S_r \gamma_1}{f_0}\right\}
  \label{eq:ch6_ndz}
\end{equation}
For $H = 2$~s, $S_r = 10$~MVA, $\gamma_1 = 0.5$~Hz/s:
NDZ corresponds to $|P_\mathrm{imb}| < 0.4$~MW (4\% of rated). The
combined criterion reduces the NDZ by additionally monitoring $V_2/V_1$
and $\Delta\theta$. Simulations show the combined NDZ is $< 1\%$ of
the $(P_g, P_l)$ operating space.

\subsection{IEC 61850 Mode-Switch Broadcast}

Upon islanding detection, the relay broadcasts a GOOSE mode-switch message
within $T_\mathrm{GOOSE} \leq 4$~ms:
\begin{equation}
  \texttt{GOOSE.mode} = \texttt{ISLAND}
  \label{eq:ch6_goose_mode}
\end{equation}
All IBRs subscribing to this GOOSE message switch from GFL to GFM mode.
The total mode-switch latency budget is:
\begin{align}
  T_\mathrm{total} &= T_\mathrm{detect} + T_\mathrm{GOOSE} + T_\mathrm{switch}
  \nonumber \\
  &\leq 60 + 4 + 50 = 114\;\mathrm{ms} < 120\;\mathrm{ms}
  \label{eq:ch6_latency}
\end{align}
where $T_\mathrm{detect} = 60$~ms (3 ROCOF cycles), $T_\mathrm{GOOSE} = 4$~ms,
and $T_\mathrm{switch} = 50$~ms (GFM activation time). This satisfies
the 120~ms target and the IEEE~1547-2018 ride-through compliance requirement
(which requires a 160~ms detection time for normal voltage conditions).

\section{Adaptive OC Setting-Group Architecture}
\label{sec:ch5_adaptiveoc}

\subsection{Motivation}

The IEC inverse-time overcurrent relay has two setting parameters:
pickup current $I_\mathrm{pickup}$ and time-multiplier setting $\TMS$.
In grid-connected mode:
\begin{equation}
  I_\mathrm{pickup}^\mathrm{grid} = 1.5\,I_\mathrm{load,max}, \quad
  \TMS^\mathrm{grid} = 0.15
  \label{eq:ch6_sg_settings}
\end{equation}
In islanded mode, the fault current is 4--8$\times$ lower. The pickup
threshold must be reduced to remain above load:
\begin{equation}
  I_\mathrm{pickup}^\mathrm{island} = 1.25\,I_\mathrm{load,max}, \quad
  \TMS^\mathrm{island} = 0.05
  \label{eq:ch6_island_settings}
\end{equation}
The transition between these settings must occur within $T_\mathrm{switch}$
to avoid both missed trips (if settings are too high) and sympathetic
trips (if settings are too low during the transition window).

\subsection{Setting-Group Pre-Arming Logic}

The setting-group switch is pre-armed when the GOOSE mode-switch signal
is received, and activates when the fault occurs or at
$t = T_\mathrm{armed} + 50$~ms (whichever comes first):
\begin{equation}
  \texttt{SG.active} = \begin{cases}
    \texttt{SG2} & \text{if } \texttt{GOOSE.mode = ISLAND} \text{ AND } t > T_\mathrm{armed} + 5\,\mathrm{ms} \\
    \texttt{SG1} & \text{otherwise (grid-connected)}
  \end{cases}
  \label{eq:ch6_sg_logic}
\end{equation}
The 5~ms delay prevents the relay from switching to the lower settings
during the transient portion of the mode-switch event (when fault current
is temporarily elevated).

\subsection{Bus Differential 87B for Microgrid Busbar}

For the microgrid common busbar, a bus differential (87B) element protects
against busbar faults that the line OC elements cannot detect:
\begin{equation}
  I_\mathrm{bus,diff} = \sum_{k} \phasor{I}_k > k_m^\mathrm{bus} \cdot I_\mathrm{bus,bias}
  \label{eq:ch6_87b}
\end{equation}
The bus differential is immune to the mode transition because it directly
measures the algebraic sum of all currents entering the busbar—a quantity
that is zero for any through condition regardless of $\kibr$.

\subsection{Coordination Table: Grid-Connected to Islanded}

Table~\ref{tab:ch6_settings} provides the relay settings for both
operating modes.

\begin{table}[htbp]
\caption{Adaptive relay settings for SAMBP microgrid (grid-connected and islanded)}
\label{tab:ch6_settings}
\begin{tabularx}{\textwidth}{@{}lXXXl@{}}
\toprule
Relay & Setting & Grid-connected & Islanded & Element \\
\midrule
R-feeder & $I_\mathrm{pickup}$ & 1.5~pu & 1.25~pu & ITOC~51 \\
R-feeder & $\TMS$ & 0.15 & 0.05 & ITOC~51 \\
R-bus & $k_m^\mathrm{bus}$ & 0.30 & 0.20 & 87B \\
R-bus & $I_{0,\mathrm{bus}}$ & 0.10~pu & 0.05~pu & 87B \\
R-gen & $I_\mathrm{pickup}^\mathrm{OV}$ & 1.10~pu & 1.20~pu & OV~59 \\
R-gen & $\TMS^\mathrm{UV}$ & 0.10~s & 0.05~s & UV~27 \\
\bottomrule
\end{tabularx}
\end{table}

\section{Microgrid Integration Study}
\label{sec:ch5_integration}

The complete microgrid protection architecture is validated through the
SAMBP microgrid integration study (SAMBP~TR-42).
The test system consists of:
\begin{itemize}
\item 4 GFL inverters (2~MW each), switchable to GFM.
\item 1 SG (5~MW), used as synchronous condenser.
\item 3 feeders with mixed load ($\cos\phi = 0.9$).
\item IEC~61850 GOOSE network (2~ms round-trip latency).
\end{itemize}


\section{Summary}
\label{sec:ch5mg_summary}

This chapter has developed the protection architecture for IBR microgrids:
\begin{enumerate}
\item \textbf{GFL/GFM characterisation:} Fault current levels are
      4--8$\times$ lower in islanded mode, necessitating dual setting groups.
      Table~\ref{tab:ch6_fault_comparison} provides quantitative data for
      setting group design.
\item \textbf{ROCOF islanding detection:} The combined ROCOF +
      vector surge + negative-sequence criterion achieves NDZ $< 1\%$
      and total isolation within 114~ms, within the 120~ms target and
      compliant with IEEE~1547-2018.
\item \textbf{Adaptive OC setting-group:} The pre-arming logic
      switches from grid-connected to islanded settings within 5~ms of
      GOOSE receipt, eliminating the mode-transition blind spot.
\item \textbf{Bus differential 87B:} Mode-invariant busbar protection
      provides 100\% coverage for busbar faults at any $\kibr$.
\end{enumerate}
